% Evidence-focused interpretation; original predictions are unchanged.
\subsection{What changes across the four settings?}
\paragraph{A: Predictable environment.}
At 65,536 interactions, PPO reaches 20\% development success against Q's 2\%
in A. B instead gives Q a small early lead, 6\% against 0\%. This does not
support a general early Q advantage. Final A success favours PPO, but its
80\% average comprises four 100\% agents and one 0\% agent, not uniformly
reliable retraining. With only ten development episodes per checkpoint,
individual curve points are noisy. The table tests development-selected
weights, which can precede the last checkpoint, while charging full-run costs.

\paragraph{B: Uncertain environment.}
PPO exceeds Q by 35.2 percentage points on average, with higher observed
success in four of five paired training seeds. Omitting any one paired seed leaves a positive mean
advantage of 29.3--44.5 points: a descriptive sensitivity check, not a
significance test. However, PPO still fails in 32.6\% of episodes.
Q's A-to-B mean decline of only 0.7 points hides paired changes from
$-71$ to $+48$ points. Thus the mean neither demonstrates robustness nor
supports our predicted larger Q loss. Aggregate outcomes cannot establish
that either agent chose safer routes.

\paragraph{C: Changed incentives.}
Relative to B, success falls by 5.3 points for Q and 14.7 for PPO. Mean
successful lengths shorten from 655 to 625 decisions and 578 to 533,
respectively, but this conditions on different sets of successful episodes.
With goal attainment prioritised over speed, the results do not justify
this penalty. Two of five seeds improve for each method, so this is not
evidence that step penalties invariably harm learning. We compare common
outcomes and raw scores: C raw scores are $1534.8\pm201.0$ for Q and
$1873.4\pm283.5$ for PPO; penalised returns are not directly comparable to B.

\paragraph{D: Transfer.}
Food-relative Q keys and shared PPO action scores retain meaning on the new
maze, aided by geometry-based feature extraction. Familiar compact keys can
still merge different situations. Every B agent's observed success declines
on D. The PPO policy selected on B development achieves 92.5\% in B but only
2\% in D, illustrating why source performance cannot certify transfer.

The original D run hit a natural map change. The revised boundary protocol,
locked before fresh evaluation of all ten B agents, uses external truncation
at map exit and takes no action on another maze. With $S$ successes and $C$
censored episodes among $N$, the completed-task success fraction for these
episodes is bounded by $[S/N,(S+C)/N]$.
The single censored Q episode gives a mean bound of 4.5--4.6\%, not a
confidence interval. Its maximum 0.1-point effect cannot explain the weak
transfer. Returns for censoring are in-map prefixes, not complete-task returns.
This is a post-inspection protocol repair; the previously inspected maze and
ROM intervention do not constitute an independent geometry-only domain test.

Within this budget, PPO is the stronger B workflow. Neither workflow's
failure rate or transfer evidence supports reliable operation where failures
matter. These findings concern the tested representations and procedures,
not universal algorithm superiority.
